\documentclass{amsart}
\usepackage{amssymb,latexsym,amsmath,amsthm,enumitem,mathrsfs,geometry}
\usepackage{fullpage}
%\usepackage{showkeys}
\usepackage{fancyhdr}
\usepackage{color}
\usepackage{stmaryrd}
\usepackage{mathrsfs}
\usepackage{hyperref}
\usepackage{lineno}
\usepackage{diagbox}
\usepackage{array}
\usepackage{tikz}
\usetikzlibrary{arrows}
\usetikzlibrary{positioning}
\usepackage{float}
\usepackage{bbm}
\usepackage{cleveref}
\usepackage{dsfont}
\usepackage{graphicx}
%\usepackage{pgfplots}
\usepackage{soul}
\input{format}

%MARGIN DELIMITERS
%\setlength{\oddsidemargin}{0cm}
%\setlength{\evensidemargin}{0cm}
%\setlength{\topmargin}{-.7cm}
%\setlength{\textwidth}{16cm}
%\setlength{\textheight}{23cm}

%These next commands put lines at top and bottom of page
\renewcommand{\headrulewidth}{0.4pt}
\renewcommand{\footrulewidth}{0.4pt}
%\setlength{\headheight}{-1cm}
%\setlength{\footskip}{8pt}
%\setlength{\headsep}{8pt}

\setlength{\headsep}{8pt}
\setlength{\footskip}{20pt}

\newcommand{\Q}{\mathbb{Q}}
\newcommand{\R}{\mathbb{R}}

\newcommand{\pfrak}{\mathfrak{p}}
\newcommand{\Cl}{\mathrm{Cl}}
\newcommand{\Norm}{\mathrm{Norm}}
\newcommand{\Ocal}{\mathcal{O}}
\newcommand{\Ascr}{\mathscr{A}}
\newcommand{\Bscr}{\mathscr{B}}
\newcommand{\Pscr}{\mathscr{P}}
\newcommand{\Escr}{\mathscr{E}}
\renewcommand{\thefootnote}{\fnsymbol{footnote}}
\newcommand{\Leg}[2]{\left( \frac{#1}{#2}\right)}
\newcommand{\Gal}{\mathrm{Gal}}
\newcommand{\Disc}{\mathrm{Disc}}
\newcommand{\Li}{\mathrm{Li}}
\newcommand{\li}{\mathrm{li}}
\newcommand{\Cond}{\mathrm{Cond}}
\newcommand{\GL}{\mathrm{GL}}
\newcommand{\1}{\mathbbm{1}}
\newtheorem{theorem}{Theorem}%[section]
\newtheorem{Theorem}{Theorem}%[section]
\newtheorem{lem}{Lemma}
\newtheorem{theoremletter}{Theorem}
\newtheorem{lemmaletter}{Lemma}
\newtheorem{corollary}{Corollary}
\newtheorem{Conj}{Conjecture}%[section]
\renewcommand*{\thetheoremletter}{\Alph{theoremletter}}
\renewcommand*{\thelemmaletter}{\Alph{lemmaletter}}
\setcounter{lemmaletter}{1}
\newtheorem*{theorem*}{Theorem}
\newtheorem*{conjecture*}{Conjecture}%[section]
\theoremstyle{remark}
\newtheorem*{remark}{Remark}
%AUTHOR COMMENTS
\definecolor{pink}{rgb}{1,.2,.6}
\definecolor{orange}{rgb}{0.7,0.3,0}
\definecolor{blue}{rgb}{.2,.6,.75}
\definecolor{green}{rgb}{.4,.7,.4}
\definecolor{purple}{RGB}{127,0,255}
%\newcommand{\author}[1]{\smargin{{\color{orange}{\sf $\blacktriangleright$ Author: #1}} }}
\newcommand{\Caroline}[1]{\noindent\reversemarginpar{{\color{orange}{Caroline: #1}} }}
\newcommand{\Sieg}[1]{\noindent\reversemarginpar{{\color{purple}{Sieg: #1}} }}
\newcommand{\Chacha}[1]{\noindent\reversemarginpar{{\color{green}{Chacha: #1}} }}
\newcommand{\Dan}[1]{\noindent\reversemarginpar{{\color{blue}{Dan: #1}} }}
\newcommand{\Question}[1]{\noindent\reversemarginpar{{\color{red}{Question: #1}} }}
\DeclareMathOperator{\EX}{\mathbb{E}}% expected value

\newcommand{\todo}[1]{{\color{cyan}{(to do: #1)}} }

\numberwithin{equation}{section}

%\allowdisplaybreaks

%%%%%%%%%%%%%%%%%%%%%%%%%%%%%%%%
%%%%%%%%%%%%%%%%%%%%%%%%%%%%%%%%

\begin{document}
\title[Unconditional Montgomery Pair Correlation]{An unconditional Montgomery Theorem for Pair Correlation of Zeros of the Riemann Zeta Function}

\author[Baluyot]{Siegfred Alan C. Baluyot}
\address{American Institute of Mathematics}
\email{sbaluyot@aimath.org}

\author[Goldston]{Daniel Alan Goldston}
\address{Department of Mathematics and Statistics, San Jose State University}
\email{daniel.goldston@sjsu.edu}

\author[Suriajaya]{Ade Irma Suriajaya}
\address{Faculty of Mathematics, Kyushu University}
\email{adeirmasuriajaya@math.kyushu-u.ac.jp}

\author[Turnage-Butterbaugh]{Caroline L. Turnage-Butterbaugh}
\address{Carleton College}
\email{cturnageb@carleton.edu}

\keywords{Riemann zeta-function, zeros, pair correlation, simple zeros, zero-density}
%\subjclass[2010]{}
%\thanks{$^{*}$ 

\dedicatory{Dedicated to Henryk Iwaniec on the occasion of his 75th birthday}

\begin{abstract}
Assuming the Riemann Hypothesis (RH), Montgomery proved a theorem concerning pair correlation of zeros of the Riemann zeta-function. One consequence of this theorem is that, assuming RH, at least $67.9\%$ of the nontrivial zeros are simple. Here we obtain an unconditional form of Montgomery's theorem and show how to apply it to prove the following result on simple zeros: Assuming all the zeros $\rho=\beta+i\gamma$ of the Riemann zeta-function such that $T^{3/8}<\gamma\le T$ satisfy $|\beta-1/2|<1/(2\log T)$,
%lie in the thin box $\{s=\sigma +it: |\sigma-1/2|<1/(2\log T),\ T^{3/8}<t\le T\}$, 
then, as $T$ tends to infinity, at least $61.7\%$ of these zeros are simple. The method of proof neither requires nor provides any information on whether any of these zeros are on or not on the critical line where $\beta=1/2$. We also obtain the same result under the weaker assumption of a strong zero-density hypothesis. 
\end{abstract}
\date{\today}

\maketitle

\section{Introduction and Statement of Results}

Let $\rho=\beta +i\gamma$ denote a nontrivial zero of the Riemann zeta-function $\zeta(s)$ with $\beta,\gamma \in \mathbb{R}$, that is, a zero satisfying $\beta>0$. 
The Riemann Hypothesis (RH) states that $\beta=1/2$ for all $\rho$. To study the pair correlation of zeros of the zeta-function, Montgomery \cite{Montgomery73} assumed RH and considered, for $x>0$ and $T\ge3$, the sum
\begin{equation}\label{MonF} \sum_{ 0<\gamma,\gamma' \le T} x^{i(\gamma-\gamma')} \frac{4}{4+(\gamma-\gamma')^2}.
\end{equation}
The first goal of this paper is to generalize Montgomery's pair correlation method so that it is unconditional. To this end, define, for $x>0$ and $T\ge3$,
\begin{equation}\label{F(x,T)}
F(x,T) := \sum_{\substack{\rho, \rho' \\ 0<\gamma,\gamma' \le T}} x^{\rho-\rho'}w(\rho-\rho'), \qquad \text{where} \qquad w(u) := \frac{4}{4 - u^2},
\end{equation}
where here and throughout the paper, zeros are counted with multiplicity. Note that if RH holds then \eqref{F(x,T)} agrees with \eqref{MonF}.

Following Montgomery, we normalize $F(x,T)$ by defining, for real $\alpha$,
\begin{equation}\label{MonF(alpha)} F(\alpha) := \left(\frac{T}{2\pi}\log T\right)^{-1} F(T^\alpha,T) = \left(\frac{T}{2\pi}\log T\right)^{-1} \sum_{\substack{\rho, \rho' \\ 0<\gamma,\gamma' \le T}} T^{\alpha(\rho-\rho')}w(\rho-\rho').\end{equation}
The first result of this paper is the following unconditional theorem.
\begin{theorem}\label{thm1}
The function $F(\alpha)$ is real, even, and nonnegative. Moreover, as
$T \to \infty$, we have
\begin{equation}\label{Fasymp} F(\alpha) =
 T^{-2\alpha}(\log T + O(1))+ \alpha + O\left(\frac1{\sqrt{\log T}}\right) 
\end{equation}
uniformly for $0\le \alpha \le 1$.
\end{theorem}
\Cref{thm1} is nearly identical to Montgomery's theorem in \cite{Montgomery73} and \cite[Lemma 8]{GM87} except it does not assume RH, and it includes the improvements from \cite[Lemma 8]{GM87} where \eqref{Fasymp} holds up to $\alpha=1$ with explicit error terms. The proof is also nearly identical. See also \cite{IK04}. A simple application of Theorem \ref{thm1} concerns the multiplicities of the zeros of $\zeta(s)$. We use a slight modification of a kernel due to Tsang \cite{Tsang3} to prove the following result.

\begin{theorem}\label{thm2} Suppose that all the zeros $\rho = \beta +i\gamma$ of the Riemann zeta-function with $T^{3/8}<\gamma\le T$ lie within the thin box
\begin{equation}\label{strip} \frac1{2} -\frac1{2\log T}< \beta < \frac1{2} + \frac1{2\log T}. \end{equation}
Then for any sufficiently large $T>0$, at least $61.7 \%$ of the nontrivial zeros are simple. \end{theorem}

\begin{remark}
The pair correlation method developed in this paper neither requires nor provides any information as to whether or not the nontrivial zeros of $\zeta(s)$ satisfy $\beta=1/2$.
%An interesting feature of the pair correlation method is that it neither requires nor provides any information on whether any of these zeros are on or not on the critical line where $\beta =1/2$.
\end{remark}

There are many results concerning the proportion of nontrivial zeros of the Riemann zeta-function that are simple.
%which lie in the critical strip $\{s\in\mathbb{C} : 0<{\rm Re}(s)<1\}$.
Pratt, Robles, Zaharescu, and Zeindler \cite{Pra20} have proved that more than $41.7\%$ of the zeros are on the critical line, and also that more than $40.7\%$ of the zeros are on the critical line and simple. Conrey, Iwaniec, and Soundararajan \cite{CIS13} have proved that more than $14/25=56\%$ of the nontrivial zeros of all Dirichlet $L$-functions are on the critical line and are simple. Going back to the case of the Riemann zeta-function, assuming RH, Montgomery \cite{Montgomery73} obtained from \Cref{thm1} that more than $2/3 = 66.\overline{6}\%$ of the zeros are simple, and soon after, Montgomery and Taylor \cite{Montgomery74} improved this to prove more than $ 67.2\%$ of the zeros are simple. Recently Chirre, Gonçalves, and de Laat \cite{CGL2020} obtained by this method $67.9\%$. By a mollifier method Conrey, Ghosh, and Gonek \cite{CGG98} showed on RH and an additional hypothesis that at least $19/27=70.3\overline{703}\%$ of the zeros are simple, and later Bui and Heath-Brown \cite{BuiHB} showed that this result holds on RH alone.

We can weaken the assumption that there are no zeros outside the box \eqref{strip} by using a strong zero-density hypothesis. Let $N(\sigma,T)$ denote the number of zeros $\rho=\beta +i\gamma$ with $\beta \ge \sigma$ and
$0<\gamma \le T$. 

\begin{theorem}\label{thm3} Assuming that 
\begin{equation}\label{Densitythm3} N(\sigma, T) = o\left( T^{2(1-\sigma)}\right) \qquad \text{for} \qquad \frac1{2} + \frac1{2\log T} \le \sigma \le \frac{25}{32}+\eta, \end{equation}
for any fixed $\eta>0$, then as $T\to \infty$, at least $61.7 \%$ of the nontrivial zeros of $\zeta(s)$ are simple. \end{theorem}

Selberg \cite{SelCollected2} made the conjecture that for all $\sigma>1/2$, we have
\begin{equation}\label{Selberg}
N(\sigma,T) = O\left( T^{1-c(\sigma-\frac12)} \frac{\log T}{\sqrt{\log \log T}}\right),
\end{equation}
where $c>0$ is some constant, and he stated that \eqref{Selberg} can often be used as a replacement for RH in the Selberg class. To see this, note that the conjecture is expected to hold for \emph{all} $\sigma>1/2$ and thus implies that almost all the nontrivial zeros are on the critical line $\{s\in\mathbb{C} : {\rm Re}(s)=1/2\}$. In a recent paper, Aryan \cite{Ary22} used this type of conjecture as a replacement of RH to obtain Montgomery's result on simple zeros. Our density conjecture \eqref{Densitythm3} implies all except $o(T)$ of the zeros are in the box \eqref{strip}, and if we extend this conjecture to all $\sigma>1/2$ we also obtain Montgomery's simple zero results. Iwaniec and Kowalski \cite[p. 249]{IK04} made a weaker Density Conjecture that $ N(\sigma,T) \ll T^{2(1-\sigma)} \log T$ for $1/2\le \sigma \le 1 $ and $T\ge 3$,
which, however, is too weak to be used in Theorem \ref{thm3}.


\section{Proof of Theorem \ref{thm1}}

Recall that if $\rho$ is a zero of $\zeta(s)$, then $1-\rho$, $\overline{\rho}$, and $1-\overline{\rho}$ are also zeros. Write 
\begin{equation*}%\label{delta}
\rho = \beta +i\gamma := 1/2 +\delta + i\gamma,
\end{equation*}
where $-1/2<\delta <1/2$. If $\delta\neq 0$ and $\gamma >0$, then there is another zero in the upper half-plane given by
\[ 1-\overline{\rho} = 1/2 -\delta + i\gamma.\]
%\[ 1-\overline{\rho} = 1 -\beta +i\gamma = 1/2 -\delta + i\gamma.\]
Therefore we may rewrite \eqref{F(x,T)} as 
\begin{equation}\label{F(x,T)2}
\begin{aligned}
%F(x,T) = \sum_{\substack{\rho, \rho' \\ 0<\gamma,\gamma' \le T}} x^{\rho-\rho'}w(\rho-\rho')
F(x,T) = \sum_{\substack{\rho, \rho' \\ 0<\gamma,\gamma' \le T}} x^{\rho+\overline{\rho'}-1}w(\rho+\overline{\rho'}-1)= \sum_{\substack{\rho, \rho' \\ 0<\gamma,\gamma' \le T}} x^{\delta +\delta' +i(\gamma-\gamma')}w(\delta +\delta' +i(\gamma-\gamma')).
\end{aligned}
\end{equation}

\begin{lem}[Montgomery]\label{lem1} Let $\rho= 1/2+\delta +i\gamma$. Then for $x\ge 1$ and all $t$ we have
\begin{equation}\label{Explicit1}
\begin{split}
\sum_\rho \frac{ 2x^{\delta + i(\gamma-t)}}{1+((t-\gamma)+i\delta)^2} = 
-\sum_{n=1}^\infty\frac{\Lambda(n)}{n^{1/2+it}} & \min\big\{\frac{n}{x},\frac{x}{n}\big\} + x^{-1} \left(\log( |t|+2) + O(1) \right) \\ & + O\left(\frac{x^{1/2} }{1+t^2}\right) + O\left(\frac{x^{-5/2} }{|t|+2}\right).
\end{split}
\end{equation}
\end{lem}
\begin{proof}[Proof of Lemma \ref{lem1}] This is the Lemma from \cite{Montgomery73} if one takes $\sigma=3/2$ and $\delta=0$. 
The starting point for proving this lemma is the explicit formula due to Landau \cite{Lan09} that, for $x>1$ and $x\neq p^m$,
\begin{equation}\label{LandauExplicit}
\sum_\rho \frac{x^{\rho-s}}{s-\rho} = \sum_{n\le x}\frac{\Lambda(n)}{n^s} + \frac{\zeta'}{\zeta}(s) - \frac{x^{1-s}}{1-s} - \sum_{n=1}^\infty \frac{x^{-2n-s}}{2n+s},
\end{equation}
provided $s\neq 1$, $s \neq \rho$, $s\neq -2n$, which we henceforth assume. When $s=0$ this is the usual explicit formula for primes. Writing $s=\sigma+it$ and $\rho=1/2+\delta+i\gamma$, we obtain
\[ \sum_{\rho=1/2+\delta+i\gamma} \frac{x^{1/2+\delta-\sigma+i(\gamma-t)}}{\sigma-1/2-\delta+i(t-\gamma)} = R(\sigma+it), \]
where $R(\sigma+it)$ is the right-hand side of \eqref{LandauExplicit} which does not depend on $\rho$ and is treated exactly as in \cite{Montgomery73}.
Multiplying both sides by $x^{\sigma-1/2}$, we obtain
\[\sum_\rho \frac{x^{\delta+i(\gamma-t)}}{\sigma-1/2-\delta+i(t-\gamma)} = x^{\sigma-1/2}R(\sigma+it).\]
Next, replace $\sigma$ with $1-\sigma$ in the equation above which adds the conditions $s\neq 0$ and $s\neq 2n+1$ and gives
\[
\sum_\rho \frac{x^{\delta+i(\gamma-t)}}{1/2-\sigma-\delta+i(t-\gamma)} = x^{1/2-\sigma}R(1-\sigma +it). \]
Subtract this equation from the previous one and simplify to obtain
\[ \sum_\rho \frac{(2\sigma-1)x^{\delta+i(\gamma-t)}}{(\sigma - \frac12)^2+((t-\gamma)+i\delta)^2} = x^{\sigma-1/2}R(\sigma+it) - x^{1/2-\sigma}R(1-\sigma +it).\]
Taking $\sigma =3/2$ gives the left-hand side of the lemma, and the right-hand side is obtained exactly as in the original proof. 
\end{proof}

\begin{lem}\label{lem2} Letting $N(T)$ denote the number of zeros in the upper half plane up to height $T$, we have
\begin{equation}\label{N(T)} N(T) := \sum_{0 < \gamma \le T} 1= \frac{T}{2\pi} \log \frac{T}{2 \pi} - \frac{T}{2\pi} + O(\log T). \end{equation}
\end{lem}
\noindent This is proved in many books, for instance \cite[Theorem 25]{Ingham1932}, \cite[Theorem 9.4]{Titchmarsh}, or \cite[Corollary 14.3]{MontgomeryVaughan2007}. 
In particular, we have 
\begin{equation}\label{zeroestimte} N(T) \sim \frac{T}{2\pi}\log T ,\qquad \text{and}\qquad N(T+1)-N(T) \ll \log T. \end{equation}

\begin{lem}\label{lem3}
We have, for $x>0$ and $T\ge 3$,
\begin{equation}\label{lem3eq} F(x,T) = \frac{2}{\pi} \int_{-\infty}^{\infty} \left|\sum_{\substack{\rho \\ 0<\gamma \le T}} \frac{x^{\rho-1/2}}{1-\left(\rho-(1/2+it)\right)^2} \right|^2 \, dt. \end{equation}
\end{lem}

\begin{proof}[Proof of Lemma \ref{lem3}] We will make use of the formula, with $a\in \mathbb{C}$ and $-1< \text{Im}(a)<1$,
\begin{equation}\label{integral}
\begin{aligned}
\int_{-\infty}^\infty \frac{dt}{(1+t^2)(1+(t+a)^2)} &= 2\pi i \left( \frac1{2i(1+(a+i)^2)} + \frac1{2i(1+(-a+i)^2)} \right)= \frac{2\pi}{4+a^2},
\end{aligned}
\end{equation}
which is easily obtained by residues or Mathematica.
Now, multiplying out the right-hand side of \eqref{lem3eq}, we obtain 
\[ \begin{split} \frac{2}{\pi} \sum_{\substack{\rho, \rho' \\ 0<\gamma,\gamma' \le T}} x^{\rho+\overline{\rho'}-1}& \int_{-\infty}^\infty \frac{dt}{\big(1- \left(\rho -(1/2+it)\right)^2\big)\big(1- \left(\overline{\rho'} -(1/2-it)\right)^2\big)}\\&
=\frac{2}{\pi} \sum_{\substack{\rho, \rho' \\ 0<\gamma,\gamma' \le T}} x^{\rho+\overline{\rho'}-1}
\int_{-\infty}^\infty \frac{dt}{\big(1+ \left(t+i(\rho -1/2)\right)^2\big)\big(1+ \left(t-i(\overline{\rho'} -1/2)\right)^2\big)} \\&
=\frac{2}{\pi} \sum_{\substack{\rho, \rho' \\ 0<\gamma,\gamma' \le T}} x^{\rho+\overline{\rho'}-1}
\int_{-\infty}^\infty \frac{dt}{\big(1+ t^2\big)\big(1+ \left(t-i(\rho +\overline{\rho'} -1 )\right)^2\big)}\\&
= \sum_{\substack{\rho, \rho' \\ 0<\gamma,\gamma' \le T}} x^{\rho+\overline{\rho'}-1}\frac{4}{4+ \big(i(\rho +\overline{\rho'} -1)\big)^2}\\& 
= F(x,T)
\end{split}\]
by \eqref{F(x,T)2}.
\end{proof}

We now rewrite \eqref{lem3eq} with $\rho = 1/2+\delta +i\gamma$ so that
\begin{equation}\label{lem3eq2} F(x,T) = \frac{2}{\pi} \int_{-\infty}^{\infty} \left|\sum_{\substack{\rho \\ 0<\gamma \le T}} \frac{x^{\delta+i\gamma}}{1+ \left((t-\gamma)+i\delta\right)^2} \right|^2 \, dt. \end{equation}
Define 
\begin{equation}\label{Theta(t)} \Theta(t) := \max\left\{\beta : \rho = \beta +i\gamma,\ 0<|\gamma|\le t\right\}. \end{equation}
If RH is false then $\Theta(t)$ is a step function, and we often replace it with the \lq \lq zero-free" region in the complex $s=\sigma +i t$ plane
\begin{equation}\label{eta-zerofree} \sigma> 1-\eta(t) \ge \Theta(t), \qquad t\ge 3,\end{equation}
where $0<\eta(t)\le 1/2$ and $\eta(t)$ is a continuous decreasing but not necessarily strictly decreasing function. Clearly we can make $\eta(t)$ as close to $1-\Theta(t)$ as we wish pointwise except at the jumps of $\Theta(t)$. 
We will often make use of this with $\rho = 1/2+\delta + i\gamma$ in the form
\begin{equation}\label{delta-Theta} \delta \le \Theta(t)-1/2 \le 1/2-\eta(t) \qquad \text{for} \qquad 0<\gamma \le t.\end{equation}

Our next lemma gives the unconditional version of Montgomery's result that relates $F(x,T)$ to the explicit formula in Lemma \ref{lem1}.

\begin{lem}\label{lem4}
For $x\ge 1$ and $T\ge 3$, let 
\begin{equation}\label{L(x,T)} L(x,T) := \int_{0}^{T} \left|\sum_\rho \frac{2x^{\delta+i\gamma}}{1+ \left((t-\gamma)+i\delta\right)^2} \right|^2 \, dt. \end{equation}
Then
\begin{equation}\label{lem4eq} F(x,T) = \frac1{2\pi} L(x,T) + O\left( x^{1-2\eta(T\log^2T)}\log^3 T\right)+O(x). \end{equation}
\end{lem}


\begin{proof}[Proof of Lemma \ref{lem4}] We first truncate the sum over zeros in \eqref{L(x,T)} using trivial estimates, which then allows us to apply Montgomery's argument without modification to the truncated sum. For $x\ge 1$, since $|\delta|< 1/2$, we have the trivial estimate
\begin{equation}\label{trivialestimate} \left|\sum_\rho \frac{2x^{\delta+i\gamma}}{1+ \left((t-\gamma)+i\delta\right)^2 }\right|\ll x^{1/2}\sum_\gamma \frac1{1+(t-\gamma)^2} \ll x^{1/2}\log(|t|+2),\end{equation}
where the last estimate is well known and follows from the second estimate in \eqref{zeroestimte}.
In the same way we have, for $0\le t \le T$ and $Z\ge 2T$,
\begin{equation}\label{trivialestimate2} \left|\sum_{\substack{ \rho \\ |\gamma| \ge Z}}\frac{2x^{\delta+i\gamma}}{1+ \left((t-\gamma)+i\delta\right)^2 }\right|\ll x^{1/2}\sum_{\gamma\ge Z} \frac1{\gamma^2} \ll \frac{x^{1/2}\log Z}{Z}.\end{equation}

Next, on squaring we have
\[ L(x,T) = 4 \sum_{\rho, \rho'} x^{\delta+\delta'+i(\gamma-\gamma')} \int_0^T \frac{dt}{(1+ \left((t-\gamma)+i\delta\right)^2)(1+ \left((t-\gamma')-i\delta'\right)^2)}. \]
By \eqref{trivialestimate} and \eqref{trivialestimate2} we can exclude the terms with $\gamma \not \in [-Z,Z]$ in the sum above with an error 
\[\begin{split} &\ll \int_0^T \left|\sum_{\substack{ \rho \\ |\gamma| \ge Z}}\frac{x^{\delta+i\gamma}}{1+((t-\gamma)+i\delta))^2}\right|\left|\sum_{ \rho'}\frac{x^{\delta'-i\gamma'}}{1+((t-\gamma')-i\delta')^2}\right|\, dt 
\ll \frac{xT\log^2Z}{Z}.\end{split} \]
Taking $Z=T \log^2 T$, we conclude that
\[ L(x,T) = 4 \sum_{\substack{\rho, \rho'\\ |\gamma|\le Z, |\gamma'|\le Z}} x^{\delta+\delta'+i(\gamma-\gamma')} \int_0^T \frac{dt}{(1+ \left((t-\gamma)+i\delta\right)^2)(1+ \left((t-\gamma')-i\delta'\right)^2)} + O(x). \]
 Montgomery, arguing unconditionally except for taking $\delta=\delta'=0$ in $L(x,T)$ and with no truncation, showed the terms with $\gamma\not \in [0,T]$ can be excluded with an error $O(\log^3 T)$, and then the range of integration can be extended to $\mathbb{R}$ with an error $O(\log^2T)$. Here we apply the same argument where we need to include the factor 
 \begin{equation}\label{xdependence} x^{\delta +\delta' +i(\gamma-\gamma')} \ll x^{2\Theta(Z)-1}\le x^{1-2\eta(T\log^2 T)}\end{equation}
in the error term, at which point the bound for these error terms is majorized by dropping the truncation at $Z$ which then exactly matches Montgomery's argument. Thus we obtain by \eqref{lem3eq2}
\[\begin{split} L(x,T) &= 4 \sum_{\substack{\rho, \rho'\\ 0\le \gamma, \gamma'\le T} } x^{\delta+\delta'+i(\gamma-\gamma')} \int_{-\infty}^{\infty} \frac{dt}{(1+ \left((t-\gamma)+i\delta\right)^2)(1+ \left((t-\gamma')-i\delta'\right)^2)} \\&
\hskip 2in +O(x^{1-2\eta(T\log^2 T)}\log^3T)+ O(x) \\&
= 2\pi F(x,T) +O(x^{1-2\eta(T\log^2 T)}\log^3T)+ O(x).
\end{split}\]
\end{proof}

\begin{proof}[Proof of Theorem \ref{thm1}]
Since $w(u)$ is even, we see from \eqref{F(x,T)} that
$F(1/x,T) = F(x,T)$, and therefore $F(\alpha)$ is even. That $F(\alpha)$ is real and nonnegative follows immediately from Lemma \ref{lem3}.
 
We write \eqref{Explicit1} as $l(x,T)=r(x,T)$, and define
\begin{equation}\label{L=R} L(x,T) := \int_0^T\left|l(x,T)\right|^2\, dt = \int_0^T\left|r(x,T)\right|^2\, dt =: R(x,T). \end{equation}
As we just saw in Lemma \ref{lem4}, 
\[ L(x,T) = 2\pi F(x,T) +O(x^{1-2\eta(T\log^2 T)}\log^3T)+ O(x).\]
The current widest known zero-free region $\sigma \ge 1-\eta(t)$ was obtained independently by Korobov and Vinogradov with
\[ \eta(t) = \frac{c}{(\log{t})^{2/3}(\log\log{t})^{1/3}}, \qquad \text{for}~ t\ge3,\] 
for some constant $c>0$. Thus we see that, for $1\le x\le T^{1/2}$,
\[x^{1-2\eta(T\log^2 T)}\log^3T \ll T^{1/2}\log^3T,\]
while for $T^{1/2}\le x \le T$
\[x^{1-2\eta(T\log^2 T)}\log^3T \ll x \exp\left(-c\frac{\log x}{(\log{T})^{2/3}(\log\log{T})^{1/3}}\right)\log^3T \ll x. \] 
We conclude for $1\le x\le T$
\begin{equation}\label{Lfinal} L(x,T) = 2\pi F(x,T) + O(T) + O(x). \end{equation}

Next, $R(x,T)$ does not depend on RH, and Montgomery\cite{Montgomery73} proved unconditionally that 
\[ R(x,T) = (1+o(1)) T x^{-2} \log^2 T + T(\log x +O(1)) +O(x\log x). \]
From \cite{GM87} this was improved, so that for $0\le x\le T$,
\begin{equation}\label{Rfinal} R(x,T) = x^{-2} T\log T(\log T +O(1)) + T(\log x + O(\sqrt{\log T})). \end{equation}
Since $L(x,T)=R(x,T)$, \eqref{Lfinal} and \eqref{Rfinal} prove Theorem \ref{thm1} on using \eqref{MonF(alpha)} to convert $F(x,T)$ to $F(\alpha)$.
\end{proof}

\begin{remark}
In \eqref{Rfinal} we have removed an extraneous factor of $\log\log T$ from \cite{GM87} which can be avoided using Lemma 6 there in place of Lemma 7. The statement there has also been corrected slightly. See also \cite{Gold81} and \cite{LPZ17}. Montgomery and Vaughan in the forthcoming book Multiplicative Number Theory II have obtained a significantly refined version of {Montgomery's theorem}.
\end{remark}


\section{Sums over Differences of Zeros}

Let $g(\alpha)\in L^{\!1}(\mathbb R)$, and, for $z=x+iy$, $x,y \in \mathbb{R}$, define the Fourier transform $\widehat{g}(z)$ of $g(\alpha)$ by
\begin{equation}\label{hat-g}
\widehat{g}(z) = \int_{-\infty}^\infty g(\alpha)e(-z\alpha)\,d\alpha, \qquad \text{where} \qquad e(w) = e^{2\pi i w}.
\end{equation}
Thus $\widehat{g}(z)$ is an analytic function for all $z$. 
Taking $z= i(\rho-\rho')\frac{\log T}{2\pi}$, we have 
\[ \widehat{g}\left( i(\rho-\rho')\frac{\log T}{2\pi}\right) = \int_{-\infty}^\infty g(\alpha)T^{\alpha(\rho -\rho')}\,d\alpha.
\]
Multiplying both sides of this equation by $w(\rho-\rho')$ and summing over $0<\gamma, \gamma' \le T$, we obtain
\begin{equation}\label{pairsum}
\sum_{\substack{\rho, \rho' \\ 0<\gamma,\gamma'\le T}} \widehat{g}\left(i(\rho -\rho')\frac{\log T}{2\pi}\right) w(\rho-\rho') 
= \left(\frac{T}{2\pi}\log T\right)\int_{-\infty}^\infty F(\alpha)g(\alpha)\,d\alpha.
\end{equation}

We now apply Theorem \ref{thm1} to obtain the following unconditional version of Montgomery's result on evaluating sums over pairs of zeros for even kernels with Fourier transforms supported in $[-1,1]$. If we assume RH this agrees with the earlier version in \cite{Montgomery73}. 

\begin{lem}\label{lem5} Suppose $ \alpha \in \mathbb{R}$ and $z\in \mathbb{C}$. Suppose $r(\alpha)$
is a real-valued even function in $L^{\!1}(\mathbb R)$ with support in $[-1,1]$, and also $r(\alpha)$ is Lipschitz continuous at $\alpha =0$. Then $\widehat r(z)$ is an even analytic function,
\begin{equation}\label{rhateven}
\widehat{r}(z) = 2\int_{0}^1 r(\alpha) \cos(2 \pi z\alpha)\,d\alpha, \end{equation}
and we have
\begin{equation}\label{lem5eq}
\sum_{\substack{\rho, \rho' \\ 0<\gamma,\gamma'\le T}} \widehat{r}\left(i(\rho -\rho')\frac{\log T}{2\pi}\right) w(\rho-\rho') = \frac{T}{2\pi}\log T
\left(r(0) + 2\int_{0}^1 \alpha r(\alpha)\,d\alpha+O\left(\frac1{\sqrt{\log T}}\right)\right).
\end{equation}
\end{lem}
Recall that a function $f(x)$ is Lipschitz continuous at a point $x=a$ if there are constants $C>0$ and $\delta >0$ such that $|f(x)-f(a)| \le C |x-a|$ for all $x$ in a neighborhood $|x-a|< \delta$ of $a$. 

\begin{proof}[Proof of Lemma \ref{lem5}] In \eqref{pairsum} we take $g(\alpha)= r(\alpha)$, and see that \eqref{rhateven}
follows from \eqref{hat-g} by the evenness of $r$. For the integral in \eqref{pairsum}, we apply \Cref{thm1} and have
\[ \begin{split} \int_{-\infty}^\infty F(\alpha)r(\alpha)\,d\alpha &= 2\int_0^1 \left(T^{-2\alpha}(\log T +O(1))+ \alpha +O\left(\frac1{\sqrt{\log T}}\right)\right)r(\alpha) \, d\alpha\\&
=2\int_0^1 T^{-2\alpha}(\log T +O(1)) r(\alpha)\, d\alpha + 2\int_0^1\alpha r(\alpha)\, d\alpha +O\left(\frac1{\sqrt{\log T}}\right). \end{split}\]
To complete the proof, we use of the Lipschitz condition on $r(\alpha)$ at $\alpha=0$ to see the first integral is
\[ \begin{split} &= 2\int_0^{\log\log T/\log T} T^{-2\alpha}(\log T +O(1))r(\alpha) \, d\alpha+O\left( \int_{\log\log T/\log T}^1 T^{-2\alpha}\log T |r(\alpha)|\, d\alpha \right) \\&
= 2\left( r(0)+ O\left(\frac{\log\log T}{\log T}\right)\right) \int_0^{\log\log T/\log T}T^{-2\alpha}(\log T +O(1))\, d\alpha + O\left(\frac1{\log T}\int_0^1|r(\alpha )|\, d\alpha\right)\\&
= 2\left( r(0)+ O\left(\frac{\log\log T}{\log T}\right)\right)\left(\frac12 +O\left(\frac1{\log T}\right)\right) + O\left(\frac1{\log T}\right) \\&
= r(0) + O\left(\frac{\log\log T}{\log T}\right).
\end{split}\]
\end{proof}


\section{Tsang's Kernel}
We define the Tsang kernel $K(z)$ through its Fourier transform by
\begin{equation}\label{TsangFT} \widehat{K}(t) := j(2\pi t){\rm sech}(2 \pi t) ,\end{equation}
where $j(\alpha)$ is an even, non-negative, bounded function supported on $|\alpha|\le 1$, and we also assume $j$ is twice differentiable on $[0,1]$ with one-sides derivatives at the endpoints. We also require, for all $w \in \mathbb{R}$,
\[ 0\le \widehat{j}(w) \ll \frac1{1+w^2}. \]

Thus the Tsang kernel $K(z)$ is
\begin{equation}\label{K(z)} K(z)= \int_{-\infty}^\infty \widehat{K}(t) e(zt)\, dt = 2\int_0^\infty j(2\pi t){\rm sech}(2\pi t)\cos(2\pi z t) \, dt 
=\frac1{\pi}\int_0^1 j(\alpha){\rm sech}( \alpha)\cos( z \alpha) \, d\alpha.\end{equation}

Tsang took the function $\widehat j$ to be the Fej\'er kernel, that is, 
\begin{equation}\label{Fejer}j_F(\alpha)=\max\{0,1-|\alpha|\}, \qquad \widehat j_F(w)= \left(\frac{\sin \pi w}{\pi w}\right)^2, \end{equation}
and we will also take $\widehat j$ to be the Montgomery-Taylor kernel \cite{Montgomery74}, \cite{CheerGold} given by
\begin{equation}\label{ M-T}
j_M(\alpha) = \frac1{1 - \cos{\sqrt{2}}}\left(\frac1{2\sqrt{2}}\sin\left(\sqrt{2}j_F(\alpha))\right) + \frac12j_F(\alpha)\cos\left(\sqrt{2}\alpha\right)\right), \end{equation}
where $j_F(\alpha)$ is given in \eqref{Fejer}, and 
\begin{equation}
\widehat j_M(w) = \frac1{1 - \cos{\sqrt{2}}}\left(\frac{\sin(\frac12(\sqrt{2}-2\pi w))}{\sqrt{2}-2\pi w} + \frac{\sin(\frac12(\sqrt{2}+2\pi w))}{\sqrt{2}+2\pi w} \right)^2.
\end{equation}


Using the properties of $j$, Tsang proved that $K$ has the following properties \cite[Lemma 1]{Tsang3}. 
\begin{lem}[K.-M. Tsang]\label{lemTsang} The kernel $K(z)$ is an even entire function such that:
\begin{enumerate}[label={(\alph*)}]
\item\label{lemTsang-a} $K(x)>0$ for all $x \in \mathbb{R}$,

\item\label{lemTsang-b} For $z\in\mathbb{C}-\{0\}$, $K(z) \ll \frac{e^{|{\rm Im}(z)|}}{|z|^2}$,

\item\label{lemTsang-c} For $z=x+iy$, $x,y\in \mathbb{R}$, then when $|y|<1$, we have ${\rm Re}\,K(x+i y)>0$.
\end{enumerate}
\end{lem}


\section{Application to sums over differences of zeros}
\label{sec-sum-diff-zeros}

Applying the Tsang kernel in Lemma \ref{lem5} we obtain the following result. 
\begin{lem}\label{lem7} We have
\begin{equation}\label{lem7a} 2\pi \sum_{\substack{\rho, \rho' \\ 0<\gamma,\gamma'\le T\\ |\beta-\beta'|<\frac1{\log T}}} 
{\rm Re}\,K\left(-i(\rho -\rho')\log{T}\right) + \mathcal{S}(T)=
\left(\widehat{K}(0) + 2\int_0^1 \alpha\widehat{K}\left(\frac{\alpha}{2\pi}\right)\, d\alpha + O\left(\frac1{\sqrt{\log T}}\right)\right)\frac{T}{2\pi}\log T,
\end{equation}
where
\begin{equation}\label{lem7b} \mathcal{S}(T) := 2\pi {\rm Re}\sum_{\substack{\rho, \rho' \\ 0<\gamma,\gamma'\le T\\ |\beta-\beta'| \ge \frac1{\log T}}} K\left(-i(\rho -\rho')\log{T}\right)w(\rho-\rho') ,
\end{equation}
and $K$ and $\widehat K$ are given in \eqref{TsangFT} and \eqref{K(z)}. Here ${\rm Re}\,K> 0$ for every term in the sum in \eqref{lem7a}, and $\widehat{K}\ge 0$.
\end{lem}
\begin{proof}[Proof of Lemma \ref{lem7}]
In Lemma \ref{lem5} we take, for $\alpha \in \mathbb{R}$,
$$ r(\alpha) = \widehat{K}\left(\frac{\alpha}{2\pi}\right) = j(\alpha){\rm sech}(\alpha), $$
so that ${\rm supp}\,r\subset[-1,1]$ and
$$ \widehat{r}(z) = 2\pi K(-2\pi z). $$
Thus \eqref{lem5eq} becomes 
\begin{equation}\label{kernel-sum}
2\pi \sum_{\substack{\rho, \rho' \\ 0<\gamma,\gamma'\le T}} 
K\left(-i(\rho -\rho')\log{T}\right) w(\rho-\rho')
= \left(\widehat{K}(0) + 2\int_0^1 \alpha\widehat{K}\left(\frac{\alpha}{2\pi}\right)\, d\alpha +O\left(\frac1{\sqrt{\log T}}\right)\right)\frac{T}{2\pi}\log T.
\end{equation} 
From Lemma \ref{lemTsang} we have
\begin{equation}\label{pos-real}
{\rm Re}\,K\left(-i(\rho -\rho')\log{T}\right)
= {\rm Re}\,K\left((\gamma-\gamma'-i(\beta-\beta'))\log T\right) > 0 \qquad \text{if} \quad |\beta-\beta'| < 1/\log T.
\end{equation}

Since $w(\rho-\rho')$ is complex-valued for zeros off the half-line, we first need to remove this weight when $|\beta-\beta'| < \frac1{\log T}$ before applying \eqref{pos-real}. This is easily done
since, noting $w(z)-1 = \frac14z^2w(z)$ and so $w(0)-1=0$, and using \ref{lemTsang-b} of Lemma \ref{lemTsang},
\begin{equation}\label{removing-weight}
\begin{aligned}
\sum_{\substack{\rho,\rho' \\ 0<\gamma,\gamma'\le T\\ |\beta-\beta'| < \frac1{\log T}}}
K\left(-i(\rho-\rho')\log{T}\right)
(w(\rho-\rho')-1) &=
\sum_{\substack{\rho\neq\rho' \\ 0<\gamma,\gamma'\le T\\ |\beta-\beta'| < \frac1{\log T}}}
K\left(-i(\rho-\rho')\log{T}\right)
(w(\rho-\rho')-1)
\\ & \ll \sum_{\substack{\rho\neq\rho' \\ 0<\gamma,\gamma'\le T\\ |\beta-\beta'| < \frac1{\log T}}} \frac{T^{|\beta-\beta'|}}{|\rho-\rho'|^2\log^2T} \frac{|\rho-\rho'|^2}{|4-(\rho-\rho')^2|} \\
&\le \frac1{\log^2T} \sum_{\substack{\rho\neq\rho' \\ 0<\gamma,\gamma'\le T}} \frac{T^{1/\log T}}{|4-(\rho-\rho')^2|} \\ &\ll \frac1{\log^2T} \sum_{0<\gamma,\gamma'\le T} \frac1{1+(\gamma-\gamma')^2} %\ll \frac1{\log^2T}(T\log^2T)
\ll T, 
\end{aligned}
\end{equation}
where the last line uses Lemma \ref{lem2} as is done for the sum over zeros in \eqref{trivialestimate}. We now remove $w(\rho-\rho')$ from the terms in \eqref{kernel-sum} with $|\beta-\beta'| < \frac1{\log T}$ and take real parts to complete the proof.
\end{proof}

\begin{remark} In applications we only need the bound in \eqref{removing-weight} to be $o(T\log{T})$, which is obtained if we sum over all pairs of nontrivial zeros $\rho,\rho'$ with $|\beta-\beta'|<\frac{(1-\epsilon)\log\log{T}}{\log{T}}$ for any $\epsilon>0$.
\end{remark}


\section{Assumptions on Zeros}

In Theorem \ref{thm2} we assume all the zeros $\rho = \beta +i\gamma$ of the Riemann zeta-function with $T^{3/8}<\gamma\le T$ lie within the thin box
\begin{equation}\label{strip2} \frac1{2} -\frac1{2\log T}< \beta < \frac1{2} + \frac1{2\log T}, \end{equation}
and in Theorem \ref{thm3} we assume the strong density hypothesis
\begin{equation}\label{Densitythm3B} N(\sigma, T) = o\left( T^{2(1-\sigma)}\right) \qquad \text{for} \qquad \frac1{2} + \frac1{2\log T} \le \sigma \le \frac{25}{32}+\eta, \end{equation}
for any fixed $\eta>0$. 

We now prove that either of these assumptions implies that,
 for any sufficiently large $T$,
\begin{equation} \label{S(T)bound} \mathcal{S}(T) = 2\pi {\rm Re}\sum_{\substack{\rho, \rho' \\ 0<\gamma,\gamma'\le T\\ |\beta-\beta'| \ge \frac1{\log T}}} K\left(-i(\rho -\rho')\log{T}\right)w(\rho-\rho') = o(T\log T). \end{equation}
To do this, we first prove that the density hypothesis \eqref{Densitythm3B} implies \eqref{S(T)bound}, and next show that the essentially stronger hypothesis \eqref{strip2} implies \eqref{Densitythm3B}, and thus \eqref{S(T)bound} by the first step. We use standard results and methods for applying zero-density results to explicit formulas, see \cite[Chapter 12]{ivic} and \cite[Chapter 10]{IK04}.

By property \ref{lemTsang-b} of Lemma \ref{lemTsang}, we have
$$ K(z) \ll \frac{e^{|{\rm Im} (z)|}} {|z|^2}, $$
and therefore
$$
K(-i(\rho -\rho')\log T)
\ll \frac{T^{|\beta-\beta'|}} {((\beta-\beta')\log T)^2+((\gamma-\gamma')\log T)^2}.
$$
Hence since $|w(\rho-\rho')|\ll 1$, we have 
\[ \mathcal{S}(T) 
\ll \sum_{\substack{\rho, \rho' \\ 0<\gamma,\gamma'\le T\\ |\beta-\beta'| \ge \frac1{\log T}}} \frac{T^{|\beta-\beta'|}} {((\beta-\beta')\log T)^2+((\gamma-\gamma')\log T)^2}.\]
By the inequality $|ab|\le \frac12(a^2 +b^2)$, we have 
\[ T^{|\beta-\beta'|} = T^{|(\beta-1/2)+(1/2-\beta')|}\le T^{|\beta-1/2|} T^{|\beta'-1/2|}\le \frac12\left(T^{|2\beta-1|}+ T^{|2\beta'-1|}\right). \]
Thus
\[ \mathcal{S}(T) 
\ll \sum_{\substack{\rho, \rho' \\ 0<\gamma,\gamma'\le T\\ |\beta-\beta'| \ge \frac1{\log T}}} \frac{T^{|2\beta-1|}+ T^{|2\beta'-1|}} {((\beta-\beta')\log T)^2+((\gamma-\gamma')\log T)^2}.\]
Since $|\beta-\beta'| \ge \frac1{\log T}$, at least one of $\beta$ or $\beta'$ must be outside the interval 
$$\left(\frac12-\frac1{2\log T},\ \frac12+\frac1{2\log T}\right). $$
By relabeling if necessary, we may assume that $\beta$ is outside this interval and that $T^{|2\beta'-1|}\leq T^{|2\beta-1|}$. Hence
\[ \begin{split} \mathcal{S}(T) &\ll \sum_{\substack{\rho=\beta+i\gamma \\ |\beta-1/2| \ge \frac1{2\log T}\\ 0<\gamma \le T}} T^{|2\beta-1|} \sum_{\substack{0<\gamma'\le T \\ |\beta-\beta'| \ge 1/\log T }} \frac1 {((\beta-\beta')\log T)^2+((\gamma-\gamma')\log T)^2}\\&
\ll \sum_{\substack{|\beta-1/2| \ge \frac1{2\log T}\\ 0<\gamma \le T}} T^{|2\beta-1|} \sum_{0<\gamma'\le T} \frac1 {1+((\gamma-\gamma')\log T)^2}. \end{split} \]
If $\rho$ is a zero of zeta then so is $1-\rho$, and $|2\beta-1|=|2(1-\beta)-1|$. Hence
\[ \mathcal{S}(T) \ll \sum_{\substack{\frac12 + \frac1{2\log T} \le \beta< 1\\ 0<\gamma \le T}} T^{|2\beta-1|} \sum_{0<\gamma'\le T} \frac1 {1+((\gamma-\gamma')\log T)^2}. \]
By \eqref{zeroestimte} we have
\begin{equation}\label{S(T)inner_sum}
\begin{aligned}
\sum_{0<\gamma'\le T} \frac1 {1 +((\gamma-\gamma')\log T)^2} &\ll \int_3^T 
\frac {\log t} {1+((t-\gamma)\log T)^2}\, dt \\&
\le \int_{-\infty}^{\infty} \frac {1} {1+ v^2}\, dv 
= \pi,
\end{aligned}
\end{equation}
and therefore we have 
\[ \mathcal{S}(T) \ll \sum_{\substack{\frac12 + \frac1{2\log T} \le \beta \le 1 \\ 0<\gamma \le T}} T^{2\beta-1}. \]
Applying Bourgain's zero-density estimate \cite{Bou00}
\begin{equation}\label{Bourgain-density}
N(\sigma,T) = o(T^{2(1-\sigma)}) \qquad\text{for}\quad 25/32+\eta\le\sigma\le1,
\end{equation} which is the hypothesis \eqref{Densitythm3B} in the remaining range. Hence
\[ \begin{split} \mathcal{S}(T) &\ll \int_{\frac12 + \frac1{2\log T}}^1 T^{2u-1}\, dN(u,T) \\&
= T^{2/\log T} N\left(\frac12 + \frac1{2\log T},T\right) - 2\log T\int_{\frac12 + \frac1{2\log T}}^1 N(u,T) T^{2u-1}\, du \\& =o(T\log T),
\end{split} \]
which proves \eqref{S(T)bound}.

We now prove that the assumption \eqref{strip2} implies \eqref{Densitythm3B}. Let $0<\eta<\frac{1}{32}$ be arbitrary and fixed, and suppose that
$$ \frac{1}{2}+\frac{1}{2\log T} \leq \sigma \leq \frac{25}{32}+\eta. $$
Then by our hypothesis we have
\begin{align*}
N(\sigma,T)
% := \#\{\rho =\beta+i\gamma : \beta\geq \sigma \text{ and } 0<\gamma\leq T\} \\
&= \#\{\rho =\beta+i\gamma : \beta\geq \sigma \text{ and } T^{3/8}<\gamma\leq T\} + \#\{\rho =\beta+i\gamma : \beta\geq \sigma \text{ and } 0<\gamma\leq T^{3/8}\} \\
&= \#\{\rho =\beta+i\gamma : \beta\geq \sigma \text{ and } 0<\gamma\leq T^{3/8}\} \\
&\leq \#\{\rho =\beta+i\gamma : 0<\gamma\leq T^{3/8}\}.
\end{align*}
The number on the last line is $N(T^{3/8})$ which is $O(T^{3/8}\log T)$ by Lemma \ref{lem2}.
Thus
\begin{align*}
N(\sigma,T) \ll T^{\frac{3}{8}+\varepsilon} 
= T^{2(1-\frac{26}{32})+\varepsilon} 
= o(T^{2(1-\sigma)})
\end{align*}
as $T\rightarrow \infty$, since
$$ 1-\frac{26}{32} < 1-\frac{25}{32}-\eta \leq 1-\sigma. $$
Hence
$$ N(\sigma,T)=o(T^{2(1-\sigma)}) \qquad\text{for}\quad \frac{1}{2}+\frac{1}{2\log T} \leq \sigma \leq \frac{25}{32}+\eta \quad\text{and fixed } 0<\eta<\frac{1}{32}, $$
as $T\rightarrow \infty$.
The estimate $N(\sigma,T)=o(T^{2(1-\sigma)})$ also holds for $\sigma\geq \frac{25}{32}+\varepsilon$, for any $\varepsilon>0$, by \cite[p. 146]{Bou00}. Therefore
$$ N(\sigma,T)=o(T^{2(1-\sigma)}) \qquad\text{for}\quad \frac{1}{2}+\frac{1}{2\log T} \leq \sigma \leq \frac{25}{32}+\eta \quad\text{and fixed } \eta>0, $$
as $T\rightarrow \infty$, which is \eqref{Densitythm3B}.


\section{Proof of Theorem \ref{thm2} and Theorem \ref{thm3}}
Letting $m_\rho$ denote the multiplicity of a zero $\rho$ of $\zeta(s)$, then
$$
\sum_{\substack{\rho \\ 0<\gamma\le T}} m_\rho = \sum_{\substack{\rho, \rho' \\ 0<\gamma,\gamma'\le T \\ \rho=\rho'}} 1
= \frac1{K(0)} \sum_{\substack{\rho=\rho' \\ 0<\gamma,\gamma'\le T}} {\rm Re}\,K\left(-i(\rho -\rho')\log{T}\right).
%= \frac1{K(0)} \sum_{\substack{\rho=\rho' \\ |\beta-\beta'|<\frac1{\log{T}} \\ 0<\gamma,\gamma'\le T}} {\rm Re}\,K\left(-i(\rho -\rho')\log{T}\right).
$$
Next note trivially that if $\rho=\rho'$, then the zeros are within the range $|\beta-\beta'|<\frac1{\log{T}}$.
By \ref{lemTsang-c} of \Cref{lemTsang}, we also have that ${\rm Re}\,K\left(-i(\rho -\rho')\log{T}\right)>0$ in the same range.
Therefore we may upper bound the sum on the right-hand side above by extending the sum to all zeros with $|\beta-\beta'|<\frac1{\log{T}}$ and obtain
\begin{equation*}%\label{sum-multiplicity-start}
\begin{aligned}
\sum_{\substack{\rho \\ 0<\gamma\le T}} m_\rho
&\le \frac1{K(0)} \sum_{\substack{\rho,\rho' \\ |\beta-\beta'|<\frac1{\log{T}} \\ 0<\gamma,\gamma'\le T}} {\rm Re}\,K\left(-i(\rho -\rho')\log{T}\right).
\end{aligned}
\end{equation*}


We proved in the last section that the assumption on zeros used in either one of Theorem \ref{thm2} or Theorem \ref{thm3} implies that
$\mathcal{S}(T) = o(T\log{T})$. 
Therefore \eqref{lem7a} of \Cref{lem7} gives
\begin{equation*}
\frac1{K(0)} \sum_{\substack{\rho,\rho' \\ |\beta-\beta'|<\frac1{\log{T}} \\ 0<\gamma,\gamma'\le T}} {\rm Re}\,K\left(-i(\rho -\rho')\log{T}\right)
\sim \frac1{2\pi K(0)}\left(\widehat{K}(0) + 2\int_0^1 \alpha\widehat{K}\left(\frac{\alpha}{2\pi}\right)\, d\alpha \right)\frac{T}{2\pi}\log T,
\end{equation*}
which implies
\begin{equation}\label{sum-multiplicity}
\sum_{\substack{\rho \\ 0<\gamma\le T}} m_\rho
\le \frac1{2\pi K(0)}\left(\widehat{K}(0) + 2\int_0^1 \alpha\widehat{K}\left(\frac{\alpha}{2\pi}\right)\, d\alpha + o(1) \right)\frac{T}{2\pi}\log T,
\end{equation}
as $T\to \infty$.
Following Montgomery's \cite{Montgomery73} argument, we see that the number of zeros which are simple satisfies
\begin{equation*}
\sum_{\substack{\rho:\,\text{simple} \\ 0<\gamma\le T}} 1 \ge \sum_{\substack{\rho \\ 0<\gamma\le T}} (2-m_\rho).
\end{equation*}
Hence, the proportion of simple zeros of $\zeta(s)$ is
\begin{equation*}
\frac1{N(T)}\sum_{\substack{\rho:\,\text{simple} \\ 0<\gamma\le T}} 1 \ge 2 - \frac1{N(T)} \sum_{\substack{\rho \\ 0<\gamma\le T}} m_\rho
\end{equation*}
which, since $N(T)\sim \frac{T}{2\pi}\log T$ by Lemma \ref{lem2} gives from \eqref{sum-multiplicity} 
\begin{equation}\label{simple-proportion}
\frac1{N(T)}\sum_{\substack{\rho:\,\text{simple} \\ 0<\gamma\le T}} 1 \ge 2 - \frac1{2\pi K(0)} \left(\widehat{K}(0) + 2\int_0^1 \alpha\widehat{K}\left(\frac{\alpha}{2\pi}\right)\, d\alpha + o(1)\right).
\end{equation}


%%%KERNEL1%%%

Suppose first we take the Fejer kernel $j(\alpha) = j_F(\alpha)$. Then $\widehat{K}(0)=1$ and computation gives
$$ 2\int_0^1 \alpha\widehat{K}\left(\frac{\alpha}{2\pi}\right)\, d\alpha
= 2\int_0^1 \alpha\left(1-\alpha\right){\rm sech}\,\alpha~ d\alpha = 0.2913876354\ldots. $$
Further, applying \eqref{K(z)}, we have upon computation that
\[ \pi K(0) = \int_0^1 j(u) {\rm sech}(u)\,du
= \int_0^1 (1-u) {\rm sech}(u)\,du = 0.4640648392\ldots. \]
Substituting these into \eqref{simple-proportion} we have
\begin{align*}
\frac1{N(T)}\sum_{\substack{\rho:\,\text{simple} \\ 0<\gamma\le T}} 1
\ge 2 - \frac{1.291387636 + o(1)}{2\times0.464064839}
= 0.608612927\ldots + o(1).
\end{align*}
Thus if all the nontrivial zeros $\rho$ of $\zeta(s)$ lie within the box $|\beta-1/2|<\frac1{2\log T}$, $T^{3/8}<\gamma \le T$, then at least $60.8\%$ of them are simple.

%%%KERNEL2%%%

We improve the above proportion to $61.7\%$ using the Montgomery-Taylor kernel $j(\alpha) =j_M(\alpha)$.
Computation gives
$$ \widehat{K}(0)= j_M(0) = 1.0061271908\ldots, $$
$$ 2\int_0^1 \alpha\widehat{K}\left(\frac{\alpha}{2\pi}\right)\, d\alpha
= 2\int_0^1 \alpha j_M(\alpha){\rm sech}\,\alpha~ d\alpha = 0.2832624869\ldots, $$
and
\[ \pi K(0) = \int_0^1 j_M(u) {\rm sech}(u)\,du = 0.4663199124\ldots. \]
Hence substituting these values into \eqref{simple-proportion} as before, we have
\begin{align*}
\frac1{N(T)}\sum_{\substack{\rho:\,\text{simple} \\ 0<\gamma\le T}} 1
\ge 2 - \frac{1.289389678 + o(1)}{2\times0.466319912}
= 0.617483786\ldots + o(1).
\end{align*}


\section*{Acknowledgement and Funding}
The authors thank the American Institute of Mathematics for its hospitality and for providing a pleasant research environment where most of the research was conducted. The first author is supported by NSF DMS-1854398 FRG. The third author was supported by JSPS KAKENHI Grant Number 22K13895. The fourth author is partially supported by NSF DMS-1902193 and NSF DMS-1854398 FRG.

\bibliographystyle{alpha}
\bibliography{AHBibliography}

\end{document}
